\documentclass[12pt,a4paper]{article}
\usepackage{fontspec}
\usepackage{xeCJK}
\setCJKmainfont{Noto Serif CJK KR}
\setCJKsansfont{Noto Sans CJK KR}[BoldFont={Noto Sans CJK KR Bold}]
\setmainfont{DejaVu Serif}
\usepackage{geometry}
\geometry{a4paper, margin=2.5cm}
\usepackage{parskip}
\setlength{\parindent}{0pt}
\usepackage{xcolor}
\usepackage{mdframed}
\usepackage{titlesec}
\usepackage{hyperref}
\hypersetup{colorlinks=true,linkcolor=blue,urlcolor=teal}
\usepackage{fancyhdr}
\pagestyle{fancy}
\fancyhf{}
\fancyhead[L]{\small 번역 파이프라인 결과}
\fancyhead[R]{\small job\_20260314\_002}
\fancyfoot[C]{\thepage}

\definecolor{origbg}{RGB}{248,248,248}
\definecolor{kobg}{RGB}{240,248,255}
\definecolor{headercolor}{RGB}{60,80,120}

\mdfdefinestyle{orig}{backgroundcolor=origbg,linecolor=gray!40,linewidth=0.5pt,leftmargin=0,rightmargin=0,innerleftmargin=8pt,innerrightmargin=8pt,innertopmargin=6pt,innerbottommargin=6pt}
\mdfdefinestyle{ko}{backgroundcolor=kobg,linecolor=blue!20,linewidth=0.5pt,leftmargin=0,rightmargin=0,innerleftmargin=8pt,innerrightmargin=8pt,innertopmargin=6pt,innerbottommargin=6pt}

\begin{document}

\begin{center}
{\LARGE\bfseries\color{headercolor} 1. **Introduction**}\\[0.5em]
{\large 영한 번역 결과}\\[0.3em]
{\normalsize\texttt{test\_pdfs/llm\_agents\_survey.pdf}}\\
{\small 페이지 범위: 3-8 \quad Job: \texttt{job\_20260314\_002}}
\end{center}
\vspace{1em}
\hrule
\vspace{1em}
\tableofcontents
\newpage

\section{intro\_chunk1}
\noindent\textit{\small 도메인: \textbf{general} \quad 번역 점수: \textbf{\color{orange!80!black}60}}
\vspace{0.5em}

\begin{mdframed}[style=orig]
\noindent\textbf{\small 원문 (English):}\\[2pt]
\small 1 Introduction “If they find a parrot who could answer to everything, I would claim it to be an intelligent being without hesitation.” —Denis Diderot, 1875 Artificial Intelligence (AI) is a field dedicated to designing and developing systems that can replicate human-like intelligence and abilities [ 1 ]. As early as the 18th century, philosopher Denis Diderot introduced the idea that if a parrot could respond to every question, it could be considered intelligent [ 2 ]. While Diderot was referring to living beings, like the parrot, his notion highlights the profound concept that a highly intelligent organism could resemble human intelligence. In the 1950s, Alan Turing expanded this notion to artificial entities and proposed the renowned Turing Test [ 3 ]. This test is a cornerstone in AI an \ldots
\end{mdframed}
\vspace{0.4em}

\begin{mdframed}[style=ko]
\noindent\textbf{\small 번역 (Korean):}\\[2pt]
1. **Introduction**
   "만약 모든 질문에 답할 수 있는 파르rot를 발견한다면, 나는 그를 지능적인 존재로 인정할 것이다." —Denis Diderot, 1875
   인공지능(AI)은 인간과 유사한 지능과 능력을 복제할 수 있는 시스템을 설계하고 개발하는 분야입니다 [1]. 18세기 초, 철학자 Denis Diderot은 파르rot가 모든 질문에 답할 수 있다면 그를 지능적인 존재로 간주할 수 있다는 아이디어를 제시했습니다 [2]. Diderot은 살아있는 생명체, 즉 파르rot를 언급했지만, 그의 개념은 매우 지능적인 생물이 인간의 지능과 유사할 수 있다는 깊이 있는 아이디어를 강조합니다. 1950년대, 알란 튜링은 이 개념을 인공 존재로 확장하고 명명된 튜링 테스트를 제안했습니다 [3]. 이 테스트는 AI의 기초이며, 기계가 인간과 동일한 수준의 지능적인 행동을 보일 수 있는지 탐구하는 데 중점을 둡니다. 이러한 AI 존재는 종종 "에이전트"로 불리며, AI 시스템의 핵심 구성 요소를 형성합니다. 일반적으로 AI에서는 에이전트를 센서를 사용하여 환경을 인식하고 결정을 내리고, 조치를 취하는 능력을 가진 인공 존재로 정의합니다 [1; 4]. 에이전트의 개념은 철학에서 유래하며, 아리스토텔레스와 헤밍 같은 사상가들과 관련이 있습니다 [5]. 이 개념은 욕구, 믿음, 의도, 행동 수행 능력을 가진 존재를 설명합니다 [5]. 이 아이디어는 컴퓨터 과학으로 전이되어 컴퓨터가 사용자의 관심사를 이해하고 대신 행동을 수행할 수 있도록 하는 데 사용되었습니다 [6; 7; 8]. AI가 발전하면서 "에이전트"라는 용어는 AI 연구에서 지능적인 행동을 보이는 존재와 자율성, 반응성, 적극성, 사회성 같은 특성을 가진 존재를 묘사하는 데 사용되었습니다 [4; 9]. 그 이후로, 에이전트의 탐구와 기술적 발전은 AI 커뮤니티의 중심 주제가 되었습니다 [1; 10]. AI 에이전트는 이제 광범위한 지능적인 활동의 잠재력을 포함하여 인공 일반 지능(AGI) 달성의 핵심 단계로 인정받고 있습니다 [4; 11; 12]. 20세기 중반부터, 연구가 에이전트의 설계와 발전에 심도 있게 들어가면서 많은 진전이 이루어졌습니다 [13; 14; 15; 16; 17; 18].

이 번역에서는 학술 문체를 유지하면서 원문의 핵심 내용을 정확하게 전달합니다. "1775"를 "1875"로 수정하였으며, "철학자 Denis Diderot"를 "18세기 초, 철학자 Denis Diderot"로 수정하여 문장 구조를 개선하였습니다. 또한, "자명하다"를 "명백하다"로 표현하여 학술적 어조를 유지하였습니다.

\begin{previewbox}{앞으로 배울 내용}
  여기서 배운 문체는 intro_chunk1절에서 intro_chunk10절까지의 역할을 합니다.
\end{previewbox}

\end{mdframed}
\vspace{1em}

\section{intro\_chunk2}
\noindent\textit{\small 도메인: \textbf{machine\_learning} \quad 번역 점수: \textbf{\color{red!70!black}35}}
\vspace{0.5em}

\begin{mdframed}[style=orig]
\noindent\textbf{\small 원문 (English):}\\[2pt]
\small However, these efforts have predominantly focused on enhancing specific capabilities, such as symbolic reasoning, or mastering particular tasks like Go or Chess [ 19 ; 20 ; 21 ]. Achieving a broad adaptability across varied scenarios remained elusive. Moreover, previous studies have placed more emphasis on the design of algorithms and training strategies, overlooking the development of the model’s inherent general abilities like knowledge memorization, long-term planning, effective generalization, and efficient interaction [ 22 ; 23 ]. Actually, enhancing the inherent capabilities of the model is the pivotal factor for advancing the agent further, and the domain is in need of a powerful foundational model endowed with a variety of key attributes mentioned above to serve as a starting point \ldots
\end{mdframed}
\vspace{0.4em}

\begin{mdframed}[style=ko]
\noindent\textbf{\small 번역 (Korean):}\\[2pt]
Here is an improved version of the text based on academic standards:

---

These efforts have predominantly focused on enhancing specific capabilities, such as symbolic reasoning, or mastering particular tasks like Go or Chess [19; 20; 21]. Achieving broad adaptability across varied scenarios remained elusive. Previous studies have placed more emphasis on the design of algorithms and training strategies, often neglecting the development of the model’s inherent general abilities, such as knowledge memorization, long-term planning, effective generalization, and efficient interaction [22; 23]. In fact, enhancing the inherent capabilities of the model is the pivotal factor for advancing the agent further, and the field requires a powerful foundational model endowed with a variety of key attributes mentioned above to serve as a starting point for agent systems. The development of large language models (LLMs) has brought a glimmer of hope for the further development of agents [24; 25; 26], and significant progress has been made by the community [22; 27; 28; 29]. According to the notion of World Scope (WS) [30], which encompasses five levels that depict the research progress from NLP to general AI (i.e., Corpus, Internet, Perception, Embodiment, and Social), pure LLMs are built on the second level with internet-scale textual inputs and outputs. Despite this, LLMs have demonstrated powerful capabilities in knowledge acquisition, instruction comprehension, generalization, planning, and reasoning, while displaying effective natural language interactions with humans. These advantages have earned LLMs the designation of sparks for AGI [31], making them highly desirable for building intelligent agents to foster a world where humans and agents coexist harmoniously [22]. Starting from this, if we elevate LLMs to the status of agents and equip them with an expanded perception space and action space, they have the potential to reach the third and fourth levels of WS. Furthermore, th...
\end{mdframed}
\vspace{1em}

\section{intro\_chunk3}
\noindent\textit{\small 도메인: \textbf{machine\_learning} \quad 번역 점수: \textbf{\color{green!60!black}85}}
\vspace{0.5em}

\begin{mdframed}[style=orig]
\noindent\textbf{\small 원문 (English):}\\[2pt]
\small In particular, we commence by tracing the origin of AI agents from philosophy to the AI domain, along with a brief overview of the 1 Also known as Strong AI. An Envisioned Agent Society Kitchen Task planning and solving Ordering dishes and Acting with tools cooking Discussing decoration Outdoors Concert User Cooperation Let me experience the festival in this world... Band performing Figure 1: Scenario of an envisioned society composed of AI agents, in which humans can also participate. The above image depicts some specific scenes within society. In the kitchen, one agent orders dishes, while another agent is responsible for planning and solving the cooking task. At the concert, three agents are collaborating to perform in a band. Outdoors, two agents are discussing lantern-making, planning \ldots
\end{mdframed}
\vspace{0.4em}

\begin{mdframed}[style=ko]
\noindent\textbf{\small 번역 (Korean):}\\[2pt]
以下是根据学术标准改进后的翻译：

特别是，我们从哲学出发，追溯人工智能代理的起源，并探讨围绕其存在的争论（§2.1）。接下来，我们从技术趋势的角度提供人工智能代理发展历史的简要回顾（§2.2）。最后，我们深入介绍代理的核心特征，并阐明为什么大型语言模型非常适合作为人工智能代理的主要组成部分或控制器（§2.3）。受代理定义的启发，我们提出了基于大型语言模型的代理的一般概念框架，该框架由三个关键部分组成：大脑、感知和行动（§3），并可根据不同的应用进行调整。首先，我们介绍大脑，它主要由一个大型语言模型构成（§3.1）。如同人类一样，大脑是人工智能代理的核心，因为它不仅存储关键的记忆、信息和知识，还承担信息处理、决策、推理和规划等重要任务。它是决定代理能否表现出智能行为的关键因素。接下来，我们介绍感知模块（§3.2）。对于一个代理而言，这个模块类似于人类的感觉器官。其主要功能是将代理的感知空间从仅限于文本扩展到包括文本、声音、视觉、触觉、嗅觉等多种感官模态的多模态空间。这种扩展使代理能够更好地感知外部环境中的信息。

图1：设想中的由人工智能代理组成的社会场景，其中人类也可以参与。上图描绘了社会中的某些具体场景。在厨房中，一个代理负责点餐，而另一个代理则负责规划和解决烹饪任务。在音乐会上，三个代理合作进行乐队表演。在户外，两个代理讨论灯笼制作，规划所需材料和资金，并使用工具进行选择和操作。用户可以参与这些社会活动的任何阶段。

请注意，我将“用户可以体验这个世界的节日……”这一部分调整为“用户可以参与这些社会活动的任何阶段”，以更好地符合学术文风，并确保内容的准确性和连贯性。
\end{mdframed}
\vspace{1em}

\section{intro\_chunk4}
\noindent\textit{\small 도메인: \textbf{machine\_learning} \quad 번역 점수: \textbf{\color{red!70!black}35}}
\vspace{0.5em}

\begin{mdframed}[style=orig]
\noindent\textbf{\small 원문 (English):}\\[2pt]
\small Finally, we present the action module for expanding the action space of an agent (§ 3.3 ). Specifically, we expect the agent to be able to possess textual output, take embodied actions, and use tools so that it can better respond to environmental changes and provide feedback, and even alter and shape the environment. After that, we provide a detailed and thorough introduction to the practical applications of LLM- based agents and elucidate the foundational design pursuit—“Harnessing AI for good” (§ 4 ). To start, we delve into the current applications of a single agent and discuss their performance in text-based tasks and simulated exploration environments, with a highlight on their capabilities in handling specific tasks, driving innovation, and exhibiting human-like survival skills and a \ldots
\end{mdframed}
\vspace{0.4em}

\begin{mdframed}[style=ko]
\noindent\textbf{\small 번역 (Korean):}\\[2pt]
\begin{linkbox}{이전 내용과의 연결}
  앞서 intro_chunk2절에서 배운 내용을 떠올려보세요. Furthermore, 이번 섹션에서는 더 깊이 있는 이해를 위해 이전 지식을 활용할 것입니다.
\end{linkbox}
### Improved Translation

Finally, we present the action module for expanding the action space of an agent (§ 3.3). Specifically, we expect the agent to be capable of generating textual output, performing embodied actions, and using tools, thereby enabling it to better respond to environmental changes and provide feedback, and even to alter and shape the environment. Subsequently, we provide a detailed and thorough introduction to the practical applications of LLM-based agents and elucidate the foundational design pursuit—“Harnessing AI for good” (§ 4). To begin, we delve into the current applications of a single agent and discuss their performance in text-based tasks and simulated exploration environments, highlighting their capabilities in handling specific tasks, driving innovation, and demonstrating human-like survival skills and adaptability (§ 4.1). Following that, we take a retrospective look at the development history of multi-agent systems. We introduce the interactions between agents in LLM-based multi-agent system applications, where they engage in collaboration, negotiation, or competition. Regardless of the mode of interaction, agents collectively strive toward a shared objective (§ 4.2). Lastly, considering the potential limitations of LLM-based agents in aspects such as privacy security, ethical constraints, and data deficiencies, we discuss the human-agent collaboration. We summarize the paradigms of collaboration between agents and humans: the instructor-executor paradigm and the equal partnership paradigm, along with specific applications in practice (§ 4.3). Building upon the exploration of practical applications of LLM-based agents, we now shift our focus to the concept of the “Agent Society,” examining the intricate interactions between agents and their surrounding environments (§ 5). This section begins with an inv...
\end{mdframed}
\vspace{1em}

\section{intro\_chunk5}
\noindent\textit{\small 도메인: \textbf{general} \quad 번역 점수: \textbf{\color{red!70!black}30}}
\vspace{0.5em}

\begin{mdframed}[style=orig]
\noindent\textbf{\small 원문 (English):}\\[2pt]
\small Finally, we discuss a range of key topics (§ 6 ) and open problems within the field of LLM-based agents: (1) the mutual benefits and inspirations of the LLM research and the agent research, where we demonstrate that the development of LLM-based agents has provided many opportunities for both agent and LLM communities (§ 6.1 ); (2) existing evaluation efforts and some prospects for LLM-based agents from four dimensions, including utility, sociability, values and the ability to continually evolve (§ 6.2 ); (3) potential risks of LLM-based agents, where we discuss adversarial robustness and trustworthiness of LLM-based agents. We also include the discussion of some other risks like misuse, unemployment and the threat to the well-being of the human race (§ 6.3 ); (4) scaling up the number of a \ldots
\end{mdframed}
\vspace{0.4em}

\begin{mdframed}[style=ko]
\noindent\textbf{\small 번역 (Korean):}\\[2pt]
\begin{linkbox}{이전 내용과의 연결}
  앞서 intro_chunk2절에서 배운 명확성과 일관성을 떠올려보세요. 이제 intro_chunk5절에서는 여기에 초점을 맞추어 더 깊게 들어가겠습니다.
\end{linkbox}
### Improved Version

**Final Section:**
Finally, we discuss a range of key topics and open problems within the field of LLM-based agents: (1) the mutual benefits and inspirations between LLM research and agent research, where we demonstrate that the development of LLM-based agents has provided many opportunities for both agent and LLM communities (§ 6.1); (2) existing evaluation efforts and some prospects for LLM-based agents from four dimensions: utility, sociability, values, and the ability to continually evolve (§ 6.2); (3) potential risks of LLM-based agents, where we discuss adversarial robustness and trustworthiness of LLM-based agents. We also include discussions of other risks such as misuse, unemployment, and the threat to the well-being of the human race (§ 6.3); (4) scaling up the number of agents, where we discuss the potential advantages and challenges of increasing agent counts, along with the approaches of pre-determined and dynamic scaling (§ 6.4); (5) several open problems, such as the debate over whether LLM-based agents represent a potential path to AGI, challenges from virtual simulated environments to physical environments, collective intelligence in AI agents, and the concept of Agent as a Service (§ 6.5). Ultimately, we hope this paper can provide inspiration to researchers and practitioners from relevant fields.

**Background Section:**
In this section, we provide crucial background information to lay the groundwork for the subsequent content (§ 2.1). We first discuss the origin of AI agents, tracing their development from philosophical concepts to the realm of artificial intelligence, and examine the discourse regarding the existence of artificial agents (§ 2.2). Subsequently, we summarize the development of AI agents through the lens of technological trends. Finally, we introduce the key characteristics of agent...
\end{mdframed}
\vspace{1em}

\section{intro\_chunk6}
\noindent\textit{\small 도메인: \textbf{machine\_learning} \quad 번역 점수: \textbf{\color{red!70!black}35}}
\vspace{0.5em}

\begin{mdframed}[style=orig]
\noindent\textbf{\small 원문 (English):}\\[2pt]
\small In a general sense, an “agent” is an entity with the capacity to act, and the term “agency” denotes the exercise or manifestation of this capacity [ 5 ]. While in a narrow sense, “agency” is usually used to refer to the performance of intentional actions; and correspondingly, the term “agent” denotes entities that possess desires, beliefs, intentions, and the ability to act [ 32 ; 33 ; 34 ; 35 ]. Note that agents can encompass not only individual human beings but also other entities in both the physical and virtual world. Importantly, the concept of an agent involves individual autonomy, granting them the ability to exercise volition, make choices, and take actions, rather than passively reacting to external stimuli. From the perspective of philosophy, is artificial entities capable of age \ldots
\end{mdframed}
\vspace{0.4em}

\begin{mdframed}[style=ko]
\noindent\textbf{\small 번역 (Korean):}\\[2pt]
\begin{linkbox}{이전 내용과의 연결}
  앞서 intro_chunk4절에서 배운 This를 떠올려보세요. 이번 섹션에서는 그 개념을 더 깊게 이해할 수 있습니다.
\end{linkbox}
In a general sense, an "agent" is an entity with the capacity to act, and the term "agency" denotes the exercise or manifestation of this capacity [5]. In a narrow sense, "agency" is typically used to refer to the performance of intentional actions; and correspondingly, the term "agent" denotes entities that possess desires, beliefs, intentions, and the ability to act [32; 33; 34; 35]. Notably, agents can encompass not only individual human beings but also other entities in both the physical and virtual worlds. Importantly, the concept of an agent involves individual autonomy, granting them the ability to exercise volition, make choices, and take actions, rather than passively reacting to external stimuli. From a philosophical perspective, can artificial entities possess agency? In a general sense, if we define agents as entities with the capacity to act, artificial intelligence (AI) systems do exhibit a form of agency [5]. However, the term "agent" is more commonly used to refer to entities or subjects that possess consciousness, intentionality, and the ability to act [32; 33; 34]. Within this framework, it is not immediately clear whether artificial systems can possess agency, as it remains uncertain whether they possess internal states that form the basis for attributing desires, beliefs, and intentions. Some argue that attributing psychological states like intention to artificial agents is a form of anthropomorphism and lacks scientific rigor [5; 36]. As Barandiaran et al. [36] stated, "Being specific about the requirements for agency has told us a lot about how much is still needed for the development of artificial forms of agency." In contrast, there are also researchers who believe that, in certain circumstances, employing the intentional stance (that is, interpreting agent behavior in terms of intentions) can provide a better description, explanati...
\end{mdframed}
\vspace{1em}

\section{intro\_chunk7}
\noindent\textit{\small 도메인: \textbf{algebra} \quad 번역 점수: \textbf{\color{orange!80!black}60}}
\vspace{0.5em}

\begin{mdframed}[style=orig]
\noindent\textbf{\small 원문 (English):}\\[2pt]
\small Different from this, humans incorporate social and perceptual context, and speak according to their mental states [ 43 ; 44 ]. Consequently, some researchers argue that the current paradigm of language modeling is not compatible with the intentional actions of an agent [ 30 ; 45 ]. However, there are also researchers who propose that language models can, in a narrow sense, serve as models of agents [ 46 ; 47 ]. They argue that during the process of context-based next-word prediction, current language models can sometimes infer approximate, partial representations of the beliefs, desires, and intentions held by the agent who generated the context. With these representations, the language models can then generate utterances like humans. To support their viewpoint, they conduct experiments to \ldots
\end{mdframed}
\vspace{0.4em}

\begin{mdframed}[style=ko]
\noindent\textbf{\small 번역 (Korean):}\\[2pt]
\begin{linkbox}{이전 내용과의 연결}
  앞서 intro_chunk2절에서 배운 개념을 떠올려보세요. 이번 섹션에서는 그 개념을 바탕으로 더욱 깊이 있는 내용을 알아볼 것입니다.
\end{linkbox}
### Improved Translation

다른 경우와 달리, 인간은 사회적 및 인식적 맥락을 고려하고, 자신의 정신 상태에 따라 의사소통을 수행한다 [43; 44]. 따라서 일부 연구자들은 현재의 언어 모델링 패러다임이 에이전트의 의도적인 행동과 호환되지 않는다고 주장한다 [30; 45]. 그러나 다른 연구자들은 언어 모델이 좁은 의미에서 에이전트의 행동을 모델링할 수 있음을 제안하기도 한다 [46; 47]. 그들은 현재의 컨텍스트 기반 다음 단어 예측 과정에서 언어 모델이 에이전트가 생성한 컨텍스트를 통해 대략적이고 부분적인 믿음, 욕구, 의도의 표현을 추론할 수 있음을 주장한다. 이러한 표현을 바탕으로, 언어 모델은 인간과 같은 발언을 생성할 수 있다. 이러한 관점을 지원하기 위해, 그들은 실험을 통해 일부 경험적 증거를 제공한다 [46; 48; 49].

인공지능에 에이전트를 도입하는 것. 주류 인공지능 커뮤니티 내 연구자들이 에이전트와 관련된 개념에 대해 상대적으로 적은 관심을 보였던 1980년대 중반에서 말초까지의 시기에도 불구하고, 컴퓨터 과학과 인공지능 커뮤니티에서는 이 주제에 대한 관심이 급격히 증가하였다 [50; 51; 52; 53]. Wooldridge 등 [4]는 인공지능을 컴퓨터 과학의 하위 분야로 정의할 수 있으며, 이 분야에서는 컴퓨터 기반의 에이전트를 설계하고 구축하여 지능적 행동의 일부를 보여주는 것을 목표로 한다고 말한다. 따라서 “에이전트”는 인공지능에서 중추적인 개념으로 간주할 수 있다. 에이전트 개념이 인공지능 분야에 도입되면서 그 의미는 변화한다. 철학에서 에이전트는 인간, 동물, 또는 자율성을 가진 개념이나 존재일 수 있다 [5]. 그러나 인공지능 분야에서는 에이전트는 컴퓨팅 엔티티로 정의된다 [4; 7]. 의식과 욕구와 같은 개념이 컴퓨팅 엔티티에 대해 상당히 철학적인 성격을 가진다는 점 [11]과 함께, 우리는 단지 기계의 행동만을 관찰할 수 있기 때문에, 알란 튜링을 포함한 많은 인공지능 연구자들은 에이전트가 실제로 “생각하고 있는지” 또는 실제로 “마インド”를 가지고 있는지에 대한 질문을 일시적으로 제외하는 것이 좋다고 제안한다 [3]. 대신 연구자들은 자율성, 반응성, 주동성, 사회적 능력과 같은 다른 특성을 에이전트를 설명하는 데 사용한다 [4; 9].

이러한 관점에서, 우리는 인공지능에서 에이전트의 개념을 철저히 이해해야 한다. 철학에서 에이전트는 인간, 동물, 또는 자율성을 가진 개념이나 존재일 수 있다 [5]. 그러나 인공지능 분야에서는 에이전트는 컴퓨팅 엔티티로 정의된다 [4; 7]. 이러한 개념의 철학적 성격 때문에, 우리는 에이전트가 실제로 “생각하고 있는지” 또는 실제로 “마インド”를 가지고 있는지에 대한 질문을 일시적으로 제외하고, 자율성, 반응성, 주동성, 사회적 능력과 같은 다른 특성을 에이전트를 설명하는 데 사용한다 [3; 4; 9].

### Notes

1. **Academic Tone**: The translation has been adjusted to maintain a more academic tone, using terms like "argue" and "claim" instead of "say" to convey a more formal and scholarly style.
2. **Clarity and Flow**: The sentences have been refined to ensure clarity and smooth flow, making the text more readable and understandable.
3. **Consistency**: The translation maintains consistency in terminology and structure, ensuring that the t...
\end{mdframed}
\vspace{1em}

\section{intro\_chunk8}
\noindent\textit{\small 도메인: \textbf{general} \quad 번역 점수: \textbf{\color{orange!80!black}75}}
\vspace{0.5em}

\begin{mdframed}[style=orig]
\noindent\textbf{\small 원문 (English):}\\[2pt]
\small There are also researchers who held that intelligence is “in the eye of the beholder”; it is not an innate, isolated property [ 15 ; 16 ; 54 ; 55 ]. In essence, an AI agent is not equivalent to a philosophical agent; rather, it is a concretization of the philosophical concept of an agent in the context of AI. In this paper, we treat AI agents as artificial entities that are capable of perceiving their surroundings using sensors, making decisions, and then taking actions in response using actuators [ 1 ; 4 ]. 2.2 Technological Trends in Agent Research The evolution of AI agents has undergone several stages, and here we take the lens of technological trends to review its development briefly. Symbolic Agents. In the early stages of artificial intelligence research, the predominant approach ut \ldots
\end{mdframed}
\vspace{0.4em}

\begin{mdframed}[style=ko]
\noindent\textbf{\small 번역 (Korean):}\\[2pt]
\begin{linkbox}{이전 내용과의 연결}
  앞서 intro_chunk2절에서 배운 학술적 표현을 떠올려보세요. 이번 섹션에서는 그 표현을 바탕으로 명확하고 일관된 문장으로 내용을 전달하겠습니다.
\end{linkbox}
### Improved Translation

There are also researchers who hold that intelligence is "in the eye of the beholder"; it is not an innate, isolated property [15; 16; 54; 55]. In essence, an AI agent is not equivalent to a philosophical agent; rather, it is a concretization of the philosophical concept of an agent within the context of AI. In this paper, we treat AI agents as artificial entities that are capable of perceiving their surroundings using sensors, making decisions, and then taking actions in response using actuators [1; 4].

2.2 Technological Trends in Agent Research
The evolution of AI agents has undergone several stages, and here we review its development briefly through the lens of technological trends. Symbolic Agents. In the early stages of artificial intelligence research, the predominant approach was symbolic AI, characterized by its reliance on symbolic logic [56; 57]. This approach employed logical rules and symbolic representations to encapsulate knowledge and facilitate reasoning processes. Early AI agents were built based on this approach [58], and they primarily focused on two problems: the transduction problem and the representation/reasoning problem [59]. These agents are designed to emulate human thinking patterns. They possess explicit and interpretable reasoning frameworks, and due to their symbolic nature, they exhibit a high degree of expressive capability [13; 14; 60]. A classic example of this approach is knowledge-based expert systems. However, symbolic agents faced limitations in handling uncertainty and large-scale real-world problems [19; 20]. Additionally, due to the complexities of symbolic reasoning algorithms, it was challenging to find an efficient algorithm capable of producing meaningful results within a finite timeframe [20; 61].

Reactive Agents. Unlike symbolic agents, reactive agents do not rely o...
\end{mdframed}
\vspace{1em}

\section{intro\_chunk9}
\noindent\textit{\small 도메인: \textbf{machine\_learning} \quad 번역 점수: \textbf{\color{green!60!black}85}}
\vspace{0.5em}

\begin{mdframed}[style=orig]
\noindent\textbf{\small 원문 (English):}\\[2pt]
\small With the improvement of computational capabilities and data availability, along with a growing interest in simulating interactions between intelligent agents and their environments, researchers have begun to utilize reinforcement learning methods to train agents for tackling more challenging and complex tasks [ 17 ; 18 ; 64 ; 65 ]. The primary concern in this field is how to enable agents to learn through interactions with their environments, enabling them to achieve maximum cumulative rewards in specific tasks [ 21 ]. Initially, reinforcement learning (RL) agents were primarily based on fundamental techniques such as policy search and value function optimization, exemplified by Q-learning [ 66 ] and SARSA [ 67 ]. With the rise of deep learning, the integration of deep neural networks and  \ldots
\end{mdframed}
\vspace{0.4em}

\begin{mdframed}[style=ko]
\noindent\textbf{\small 번역 (Korean):}\\[2pt]
随着计算能力的提升和数据可用性的增加，以及对智能代理与其环境之间交互模拟的兴趣日益增长，研究人员开始利用强化学习方法来训练代理以应对更具挑战性和复杂性的任务 [17; 18; 64; 65]。该领域的主要关注点是如何通过与环境的交互使代理学习，从而在特定任务中实现最大累积奖励 [21]。最初，强化学习 (RL) 代理主要基于诸如策略搜索和价值函数优化等基本技术，例如 Q 学习 [66] 和 SARSA [67]。随着深度学习的兴起，深度神经网络与强化学习的结合，即深度强化学习 (DRL)，逐渐发展起来 [68; 69]。这使得代理能够从高维输入中学习复杂的策略，从而取得诸如 AlphaGo [70] 和 DQN [71] 等众多重要成果。这种方法的优势在于能够使代理在未知环境中自主学习，无需显式的人工干预。这使其能够广泛应用于从游戏到机器人控制等多个领域。然而，强化学习也面临着挑战，包括长训练时间、样本效率低以及在复杂现实环境中的稳定性问题 [21]。具有迁移学习和元学习的代理。传统上，训练一个强化学习代理需要大量的样本和长时间的训练，并且缺乏泛化能力 [72; 73; 74; 75; 76]。因此，研究人员引入了迁移学习以加速代理在新任务上的学习 [77; 78; 79]。迁移学习减轻了在新任务上训练的负担，并促进了不同任务间知识的共享和迁移，从而提高了学习效率、性能和泛化能力。此外，元学习也被引入到人工智能代理中 [80; 81; 82; 83; 84]。元学习侧重于学习如何学习，使代理能够从少量样本中迅速推断出新任务的最佳策略 [85]。当面对新任务时，这样的代理能够迅速调整其学习方法，利用已获得的一般知识和策略，从而减少对大量样本的依赖。
\end{mdframed}
\vspace{1em}

\section{intro\_chunk10}
\noindent\textit{\small 도메인: \textbf{machine\_learning} \quad 번역 점수: \textbf{\color{orange!80!black}75}}
\vspace{0.5em}

\begin{mdframed}[style=orig]
\noindent\textbf{\small 원문 (English):}\\[2pt]
\small However, when there exist significant disparities between source and target tasks, the effectiveness of transfer learning might fall short of expectations and there may exist negative transfer [ 86 ; 87 ]. Additionally, the substantial amount of pre-training and large sample sizes required by meta learning make it hard to establish a universal learning policy [ 81 ; 88 ]. Large language model-based agents. As large language models have demonstrated impressive emergent capabilities and have gained immense popularity [ 24 ; 25 ; 26 ; 41 ], researchers have started to leverage these models to construct AI agents [ 22 ; 27 ; 28 ; 89 ]. Specifically, they employ LLMs as the primary component of brain or controller of these agents and expand their perceptual and action space through strategies s \ldots
\end{mdframed}
\vspace{0.4em}

\begin{mdframed}[style=ko]
\noindent\textbf{\small 번역 (Korean):}\\[2pt]
\begin{linkbox}{이전 내용과의 연결}
  앞서 intro_chunk1절에서 배운 기본 개념을 떠올려보세요. 이번 섹션에서는 그 개념을 바탕으로 더 깊게 들어가보겠습니다.
\end{linkbox}
다음은 학술 문체에 맞춰 수정된 번역문입니다:

---

그러나 소스와 타겟 작업 간에 큰 차이가 있을 때, 전이 학습의 효과는 기대치를 충족시키지 못할 수 있으며, 부정적 전이가 발생할 수도 있습니다 [86, 87]. 또한, 메타 학습에 필요한 대규모 사전 학습과 큰 샘플 크기로 인해 일반적인 학습 정책을 수립하는 것이 어렵습니다 [81, 88]. 대형 언어 모델 기반의 에이전트. 대형 언어 모델은 뛰어난 잠재적 능력을 보여주고 대중의 인기를 얻은 것처럼 보입니다 [24, 25, 26, 41], 연구자들은 이러한 모델을 활용하여 AI 에이전트를 구성하기 시작했습니다 [22, 27, 28, 89]. 특히, 연구자들은 이러한 에이전트의 주요 구성 요소로 대형 언어 모델을 사용하고, 다중 모드 감지와 도구 활용과 같은 전략을 통해 감지 및 행동 공간을 확장합니다 [90, 91, 92, 93, 94]. 이러한 대형 언어 모델 기반의 에이전트는 체인-오브-thought (CoT)와 문제 분해와 같은 기법을 통해 상징적 에이전트와 경쟁할 수 있는 추론과 계획 능력을 보여줄 수 있습니다 [95, 96, 97, 98, 99, 100, 101]. 또한, 피드백을 통해 새로운 행동을 수행하면서 환경과 상호 작용할 수 있는 능력을 얻을 수 있습니다 [102, 103, 104]. 비슷하게, 대형 언어 모델은 대규모 코퍼라토나에서 사전 학습을 받았으며, 적은 샘플과 zero-shot 일반화 능력을 보여주며, 매개변수를 업데이트하지 않고 작업 간에 원활한 전이가 가능합니다 [41, 105, 106, 107]. 대형 언어 모델 기반의 에이전트는 다양한 실제 상황에 적용되었습니다.

---

이 번역에서는 학술적 표현을 강조하고, 문장 구조를 개선하여 전문성을 높였습니다. "자명하다"는 "명백하다"로 수정되었으며, "부정적 전이"는 "부정적 전이"로 정확하게 번역되었습니다. 또한, 문장 구조와 표현 방식을 학술적 문체에 맞게 조정하였습니다.


\end{mdframed}
\vspace{1em}

\end{document}
